\documentclass[11pt]{article}

\usepackage[margin=1.1in]{geometry}
\usepackage{amsmath,amssymb,amsthm,mathtools}
\usepackage{xcolor}
\usepackage[colorlinks=true,linkcolor=blue!60!black,citecolor=blue!60!black,urlcolor=blue!60!black]{hyperref}
\usepackage{enumitem}
\usepackage{algorithm}
\usepackage{algpseudocode}

\newtheorem{theorem}{Theorem}[section]
\newtheorem{proposition}[theorem]{Proposition}
\newtheorem{lemma}[theorem]{Lemma}
\newtheorem{corollary}[theorem]{Corollary}
\theoremstyle{definition}
\newtheorem{assumption}[theorem]{Assumption}
\theoremstyle{remark}
\newtheorem{remark}[theorem]{Remark}

\newcommand{\R}{\mathbb{R}}
\newcommand{\E}{\mathbb{E}}
\newcommand{\Q}{\mathbb{Q}}
\newcommand{\Pmeas}{\mathbb{P}}
\newcommand{\F}{\mathcal{F}}
\newcommand{\tr}{\operatorname{tr}}
\newcommand{\grad}{\nabla}
\newcommand{\uhat}{\hat{u}}
\DeclarePairedDelimiter{\norm}{\lVert}{\rVert}
\DeclarePairedDelimiter{\abs}{\lvert}{\rvert}

\title{Tangent-Process Regression Targets for Martingale Extensions\\
and Logarithmic Transforms of Diffusions}
\author{Notes prepared with Claude}
\date{June 10, 2026}

\begin{document}
\maketitle

\begin{abstract}
Given a diffusion $dX_s = f(X_s,s)\,ds + \sigma(X_s,s)\,dW_s$ on $[0,1]$ and a terminal
function $h$, we derive simulation-based regression targets for the martingale extension
$h(x,t) = \E[h(X_1)\mid X_t=x]$ \emph{and its spatial gradient}, using the first-variation
(tangent) process of the flow. A single trajectory, together with one reverse-mode sweep
through the unrolled simulator, produces unbiased targets for both the value and the full
$n$-dimensional gradient at every time step, in contrast to standard Euler schemes for the
associated BSDE which use only scalar value information. We then extend the construction to
the logarithmic transform $v(x,t) = \log \E[e^{r(X_1)}\mid X_t = x]$, whose BSDE has a
quadratic driver. The key structural facts are: (i) under the Doob $h$-transform the driver
vanishes and the gradient admits a driver-free tangent-process representation along tilted
paths; (ii) the tilted dynamics can be realized by a shifted-noise rollout with a
\emph{detached} control, so that reverse-mode differentiation again yields the exact discrete
targets; (iii) the importance weights correcting for an approximate control are
\emph{pathwise} (not merely in expectation) equal to one at the fixed point; and (iv) the
value targets are biased from below by exactly half the squared $L^2$ control error, which is
the Bou\'e--Dupuis variational principle. Proofs and references are given.
\end{abstract}

\tableofcontents

\section{Setting and assumptions}

Let $(\Omega,\F,(\F_s)_{s\in[0,1]},\Pmeas)$ carry an $m$-dimensional Brownian motion $W$, and
consider the It\^o SDE in $\R^n$
\begin{equation}\label{eq:sde}
dX_s = f(X_s,s)\,ds + \sigma(X_s,s)\,dW_s, \qquad s\in[t,1],\qquad X_t = x,
\end{equation}
with solution denoted $X^{t,x}_s$. We write $\sigma_k(\cdot,s)\in\R^n$ for the $k$-th column
of $\sigma\in\R^{n\times m}$, $a = \sigma\sigma^\top$, and $\grad f \in \R^{n\times n}$ for
the spatial Jacobian, $(\grad f)_{jk} = \partial_{x_k} f_j$.

\begin{assumption}\label{ass:reg}
$f$ and $\sigma$ are continuous in $(x,s)$, continuously differentiable in $x$ with bounded
spatial derivatives, and of linear growth. The terminal datum ($h$ in
Section~\ref{sec:zero}, $r$ in Section~\ref{sec:log}) is $C^1$ with $h,\grad h$ (resp.\
$r,\grad r$) of polynomial growth; in Section~\ref{sec:log} we additionally assume $r$
bounded above (so that $\E[e^{r(X_1)}]<\infty$; boundedness of $r$ suffices for the quadratic
BSDE theory we cite, and \cite{BriandHu} relaxes this to exponential moments).
\end{assumption}

Under Assumption~\ref{ass:reg} the map $x\mapsto X^{t,x}_s$ is, for a.e.\ $\omega$, a $C^1$
diffeomorphism onto its image (Kunita \cite[Thm.~4.6.5]{Kunita}), and the Jacobian
\emph{tangent process} $J_{t\to s} \coloneqq \partial X^{t,x}_s/\partial x \in \R^{n\times n}$
satisfies the variational SDE
\begin{equation}\label{eq:tangent}
dJ_{t\to s} = \grad f(X_s,s)\,J_{t\to s}\,ds
+ \sum_{k=1}^m \grad\sigma_k(X_s,s)\,J_{t\to s}\,dW^k_s,
\qquad J_{t\to t} = I_n .
\end{equation}
By Gronwall's lemma together with the Burkholder--Davis--Gundy inequality, boundedness of
$\grad f,\grad\sigma_k$ yields
\begin{equation}\label{eq:Jmoments}
\sup_{x}\;\E\Big[\sup_{s\in[t,1]}\norm{J_{t\to s}}^p\Big] < \infty
\qquad\text{for every } p\ge 1,
\end{equation}
with a constant depending only on $p$ and the Lipschitz bounds (see
\cite[Sec.~4.6]{Kunita}).

\section{The zero-driver case: martingale extension of $h$}\label{sec:zero}

Define
\begin{equation}
h(x,t) \coloneqq \E\big[h(X^{t,x}_1)\big],
\end{equation}
so that $h(X_s,s)$ is a martingale and $h(\cdot,t)$ solves the backward Kolmogorov equation
$\partial_t h + f\cdot\grad h + \tfrac12\tr(a\grad^2 h) = 0$ with $h(\cdot,1)=h$.

\subsection{Pathwise gradient representation}

\begin{theorem}[Pathwise differentiation of the martingale extension]\label{thm:pathwise}
Under Assumption~\ref{ass:reg}, $h(\cdot,t)\in C^1(\R^n)$ and
\begin{equation}\label{eq:gradrep}
\grad_x h(x,t) \;=\; \E\big[\, J_{t\to 1}^\top \,\grad h(X^{t,x}_1)\,\big].
\end{equation}
\end{theorem}

\begin{proof}
Fix $t,x$ and a direction $e\in\R^n$, $\norm{e}=1$. For $\varepsilon>0$, the mean value
theorem applied $\omega$-wise along the $C^1$ flow gives
\[
\frac{h(X^{t,x+\varepsilon e}_1) - h(X^{t,x}_1)}{\varepsilon}
= \int_0^1 \grad h\big(X^{t,x+\theta\varepsilon e}_1\big)^\top
\,\partial_x X^{t,x+\theta\varepsilon e}_1\, e \; d\theta .
\]
As $\varepsilon\downarrow 0$ the integrand converges a.s.\ to
$\grad h(X^{t,x}_1)^\top J_{t\to1}\,e$ by continuity of the flow and of its derivative in the
initial condition \cite[Thm.~4.6.5]{Kunita}. Domination: by polynomial growth of $\grad h$,
the standard moment bound $\sup_{|y-x|\le 1}\E[\sup_s |X^{t,y}_s|^p]<\infty$, and
\eqref{eq:Jmoments}, the family of integrands is bounded in $L^q(\Pmeas\otimes d\theta)$ for
some $q>1$ (Cauchy--Schwarz), hence uniformly integrable. Vitali's convergence theorem lets
us pass to the limit inside the expectation, giving the directional derivative
$\partial_e h(x,t) = \E[\grad h(X_1)^\top J_{t\to1} e]$. Continuity of the right-hand side in
$x$ (same uniform integrability plus a.s.\ continuity of $x\mapsto (X^{t,x}_1, J_{t\to1})$)
upgrades directional differentiability to $C^1$, and \eqref{eq:gradrep} follows by
transposition. For the BSDE phrasing of this representation theorem, including
non-Markovian terminal data, see Ma and Zhang \cite{MaZhang}.
\end{proof}

\subsection{BSDE and Malliavin interpretation}

Setting $Y_s = h(X_s,s)$ and $Z_s = \sigma(X_s,s)^\top \grad h(X_s,s)$, It\^o's formula and
the Kolmogorov equation give the (zero-driver) BSDE
\begin{equation}\label{eq:bsde0}
dY_s = Z_s^\top dW_s,\qquad Y_1 = h(X_1),
\end{equation}
the linear special case of Pardoux--Peng \cite{PardouxPeng}; see also
\cite{ElKarouiPengQuenez}. The classical representation $Z_t = \E[D_t Y_1 \mid \F_t]$
(Clark--Ocone / \cite{ElKarouiPengQuenez,MaZhang}), combined with the chain rule
$D_t h(X_1) = \grad h(X_1)^\top D_t X_1$ and the flow identity
$D_t X_1 = J_{t\to 1}\,\sigma(X_t,t)$, recovers \eqref{eq:gradrep} in the form
$Z_t = \sigma^\top \grad h(X_t,t)$. The point of view that the pair $(Y,Z)$ jointly satisfies
a vector BSDE with terminal condition $(h(X_1), \sigma^\top\grad h(X_1))$ --- so that the
\emph{known} terminal gradient is data, not a by-product --- is developed as the One-Step
Malliavin (OSM) scheme by Negyesi, Andersson and Oosterlee \cite{OSM}.

\subsection{Discrete targets by reverse-mode differentiation}\label{sec:discrete0}

Fix a grid $t = t_0 < t_1 < \dots < t_N = 1$ and the Euler--Maruyama recursion
\begin{equation}\label{eq:euler}
\widehat X_{i+1} = \widehat X_i + f(\widehat X_i,t_i)\,\Delta t_i
+ \sigma(\widehat X_i,t_i)\,\Delta W_i,
\qquad \Delta W_i \sim \mathcal N(0,\Delta t_i I_m)\ \text{i.i.d.}
\end{equation}
Freezing the increments $(\Delta W_i)_i$, the map $\widehat X_i \mapsto \widehat X_N$ is a
composition of smooth maps with step Jacobians
\begin{equation}\label{eq:discreteJ}
\widehat J_{i+1\leftarrow i}
= I + \grad f(\widehat X_i,t_i)\,\Delta t_i
+ \sum_k \grad\sigma_k(\widehat X_i,t_i)\,\Delta W^k_i ,
\end{equation}
which is exactly the Euler discretization of \eqref{eq:tangent}.

\begin{proposition}[Backprop computes the discrete tangent targets]\label{prop:vjp}
Let $g_i \in \R^n$ denote the gradient delivered to the intermediate node $\widehat X_i$ by
reverse-mode automatic differentiation of the scalar $h(\widehat X_N)$ through the unrolled
recursion \eqref{eq:euler} with $(\Delta W_i)$ held fixed. Then
\[
g_i = \widehat J_{N\leftarrow i}^{\,\top}\,\grad h(\widehat X_N),
\qquad
\widehat J_{N\leftarrow i} \coloneqq \widehat J_{N\leftarrow N-1}\cdots
\widehat J_{i+1\leftarrow i},
\]
i.e.\ the exact gradient of the \emph{discrete} value function
$\widehat h_i(x) \coloneqq \E[h(\widehat X_N)\mid \widehat X_i = x]$ is
$\grad \widehat h_i(x) = \E[g_i \mid \widehat X_i = x]$. Moreover one reverse sweep computes
$(g_0,\dots,g_{N-1})$ simultaneously, at a cost of $O(1)$ rollouts.
\end{proposition}

\begin{proof}
The first identity is the chain rule: reverse-mode AD computes vector--Jacobian products, and
the Jacobian of the frozen-noise composition $\widehat X_i \mapsto \widehat X_N$ is
$\widehat J_{N\leftarrow i}$ by \eqref{eq:discreteJ} and the product rule. The second
identity is Theorem~\ref{thm:pathwise} applied verbatim to the discrete-time Markov chain:
differentiation under the expectation is justified by the same uniform-integrability
argument, the moment bound \eqref{eq:Jmoments} holding for products of the matrices
\eqref{eq:discreteJ} with constants uniform in $N$ (discrete Gronwall). The cost claim is the
cheap-gradient principle of reverse-mode AD \cite{GriewankWalther}: the backward sweep visits
each node once and deposits its adjoint there.
\end{proof}

\begin{remark}[Consistency]
The targets $g_i$ are the \emph{exact} gradients of the discrete flow that the scheme itself
simulates, not $O(\Delta t)$ approximations of the continuous tangent. The regression
therefore learns the gradient of the same discrete value function whose values it learns;
discretization error enters only through $\widehat h \approx h$, with the usual weak-order-1
rate.
\end{remark}

\subsection{Sobolev regression}\label{sec:sobolev}

Since, conditionally on $\widehat X_i$, the pair $\big(h(\widehat X_N),\,g_i\big)$ is
unbiased for $\big(\widehat h(\widehat X_i,t_i),\,\grad \widehat h(\widehat X_i,t_i)\big)$,
the population loss
\begin{equation}\label{eq:sobolevloss}
\mathcal L(\theta) = \E\sum_{i=0}^{N-1}
\Big[\big(h_\theta(\widehat X_i,t_i) - h(\widehat X_N)\big)^2
+ \lambda\,\norm{\grad_x h_\theta(\widehat X_i,t_i) - g_i}^2\Big]
\end{equation}
is minimized, for every $\lambda>0$, by $h_\theta = \widehat h$ (each term is a conditional
$L^2$ projection with the same minimizer; $\lambda$ trades off variance and conditioning of
the optimization, not bias). Each trajectory supplies $(n+1)N$ scalar targets at the cost of
one rollout and one reverse sweep. This is Sobolev training; in the financial setting with
regression at $t_0$ only it is the ``differential machine learning'' of Huge and Savine
\cite{HugeSavine}, and the all-times version coincides with the driver-free specialization of
\cite{OSM}.

\begin{remark}[Variance; nonsmooth $h$]\label{rem:variance0}
The targets $J^\top\grad h(X_1)$ may be heavy-tailed when the flow is expansive
($\norm{\grad f}$ large, multiplicative noise); antithetic increments and the iterated
control variate of \cite{Vidales} (subtract
$\sum_i \grad h_\theta(\widehat X_i)^\top \sigma\,\Delta W_i$ from the value target, which by
\eqref{eq:bsde0} is exact at convergence) mitigate this. If $\grad h$ is unavailable, the
Bismut--Elworthy--Li formula \cite{ElworthyLi}
$\grad h(x,t) = \E\big[h(X_1)\,\tfrac1{1-t}\int_t^1 (\sigma^{-1}(X_s,s)J_{t\to s})^\top
dW_s\big]$ (for elliptic $\sigma$) trades the terminal gradient for an integration-by-parts
weight, at the cost of variance of order $(1-t)^{-1}$.
\end{remark}

\section{The logarithmic transform: quadratic driver}\label{sec:log}

Now let
\begin{equation}
v(x,t) \coloneqq \log \E\big[e^{r(X^{t,x}_1)}\big],
\qquad u \coloneqq e^{v},
\end{equation}
so $u(X_s,s)$ is a positive martingale, $u$ solves the Kolmogorov equation, and the
logarithmic (Hopf--Cole / Fleming) transform \cite[Ch.~VI]{FlemingSoner} yields the
Hamilton--Jacobi--Bellman equation
\begin{equation}\label{eq:hjb}
\partial_t v + f\cdot\grad v + \tfrac12 \tr\big(a\,\grad^2 v\big)
+ \tfrac12\,\norm{\sigma^\top \grad v}^2 = 0,
\qquad v(\cdot,1) = r .
\end{equation}
With $Y_s = v(X_s,s)$, $Z_s = \sigma(X_s,s)^\top\grad v(X_s,s)$, It\^o's formula gives the
quadratic BSDE
\begin{equation}\label{eq:qbsde}
dY_s = -\tfrac12\norm{Z_s}^2\,ds + Z_s^\top dW_s,\qquad Y_1 = r(X_1),
\end{equation}
well-posed for bounded $r$ with $Z\in\mathrm{BMO}$ by Kobylanski \cite{Kobylanski} (unbounded
$r$ with exponential moments: \cite{BriandHu}).

Applying Section~\ref{sec:zero} directly to $u=e^v$ is correct
($\grad u(x,t) = \E[e^{r(X_1)} J_{t\to1}^\top \grad r(X_1)]$) but suffers exponential-in-$r$
relative variance, the log-partition-function pathology. The remedy is to absorb the
exponential weight into the dynamics.

\subsection{The Doob $h$-transform removes the driver}

Since $u(X_s,s)$ is a positive martingale, It\^o on $\log u$ gives the Dol\'eans exponential
representation
\begin{equation}\label{eq:doleans}
\frac{u(X_s,s)}{u(x,t)}
= \exp\Big(\int_t^s Z_\rho^\top dW_\rho - \tfrac12\int_t^s \norm{Z_\rho}^2 d\rho\Big),
\qquad s\in[t,1].
\end{equation}
Define, conditionally on $X_t = x$, the tilted measure
\begin{equation}\label{eq:Q}
\frac{d\Q}{d\Pmeas}\Big|_{\F_1} = e^{\,r(X_1) - v(x,t)} = \frac{u(X_1,1)}{u(x,t)} .
\end{equation}
By \eqref{eq:doleans} and Girsanov's theorem ($Z\in$ BMO guarantees the Dol\'eans exponential
is a true martingale \cite{Kazamaki}), $W^\Q_s \coloneqq W_s - \int_t^s Z_\rho\,d\rho$ is a
$\Q$-Brownian motion and the $h$-transformed dynamics are
\begin{equation}\label{eq:tilted}
dX_s = \big(f + a\,\grad v\big)(X_s,s)\,ds + \sigma(X_s,s)\,dW^\Q_s .
\end{equation}
This is the classical equivalence between the logarithmic transform and entropy-regularized
stochastic control \cite{FlemingSoner,DaiPra}.

\begin{theorem}[Driver-free gradient representation under the tilt]\label{thm:tiltgrad}
Under Assumption~\ref{ass:reg},
\begin{equation}\label{eq:tiltgradrep}
\grad v(x,t) \;=\; \E_{\Q}\big[\, J_{t\to1}^\top\, \grad r(X_1)\,\big|\, X_t = x\big],
\end{equation}
where $J_{t\to s}$ is the tangent process \eqref{eq:tangent} of the \emph{uncontrolled} flow,
i.e.\ the solution of
$dJ = \grad f\,J\,ds + \sum_k \grad\sigma_k\,J\,dW^k_s$
driven by the original increments $dW_s = dW^\Q_s + Z_s\,ds$, evaluated along the tilted path
\eqref{eq:tilted}.
\end{theorem}

\begin{proof}[Proof 1 (change of measure)]
Apply Theorem~\ref{thm:pathwise} to the terminal function $e^{r}$ (polynomial growth of
$e^r\,\grad r$ holds since $r$ is bounded above and $\grad r$ has polynomial growth):
$\grad u(x,t) = \E_\Pmeas\big[e^{r(X_1)}\,J_{t\to1}^\top \grad r(X_1)\big]$. Divide by
$u(x,t)$ and use $\grad v = \grad u / u$ together with \eqref{eq:Q}:
\[
\grad v(x,t)
= \E_\Pmeas\Big[e^{\,r(X_1)-v(x,t)}\,J_{t\to1}^\top\grad r(X_1)\Big]
= \E_\Q\big[J_{t\to1}^\top \grad r(X_1)\big].
\]
The pair $(X,J)$ is a path functional of $(x, W)$; rewriting its driving noise as
$dW = dW^\Q + Z\,dt$ expresses the same functional in terms of the $\Q$-Brownian motion,
which is the stated form.
\end{proof}

\begin{proof}[Proof 2 (PDE / martingale argument; $\sigma$ constant for clarity)]
Let $p \coloneqq \grad v$. Differentiating \eqref{eq:hjb} in $x_k$ and using symmetry of the
Hessian, $\partial_k p_j = \partial_j p_k$, so that
$p^\top a\,\partial_k p = (a p)\cdot \grad p_k$:
\begin{equation}\label{eq:gradpde}
\partial_t p_k + \big(f + a p\big)\cdot \grad p_k + \tfrac12\tr\big(a\grad^2 p_k\big)
= -\big((\grad f)^\top p\big)_k .
\end{equation}
The quadratic term has recombined with $f$ into the \emph{controlled} drift of
\eqref{eq:tilted}: in BSDE language, a driver whose $z$-gradient is
$\grad_z(\tfrac12\norm{z}^2) = z$ is absorbed exactly by the Girsanov shift with kernel $Z$,
which is the $h$-transform. Now set $M_s \coloneqq J_{t\to s}^\top\, p(X_s,s)$ along the
tilted dynamics. With $\sigma$ constant, $dJ = \grad f(X_s,s)\,J\,ds$ under either measure,
and It\^o's formula under $\Q$ together with \eqref{eq:gradpde} gives
\[
dM_s = J^\top (\grad f)^\top p \, ds
+ J^\top\Big(\underbrace{\partial_t p + (f+ap)\cdot\grad p
+ \tfrac12\tr(a\grad^2 p)}_{=-(\grad f)^\top p}\Big)\,ds
+ J^\top (\grad p)^\top \sigma\, dW^\Q_s
= J^\top (\grad p)^\top\sigma\,dW^\Q_s,
\]
a local martingale; the moment bounds \eqref{eq:Jmoments} (valid under $\Q$ since
$f + a\grad v$ inherits local Lipschitz bounds on compacts, with a standard localization)
make it a true martingale on $[t,1]$, whence
$p(x,t) = M_t = \E_\Q[M_1] = \E_\Q[J_{t\to1}^\top \grad r(X_1)]$. For state-dependent
$\sigma$, Proof~1 applies unchanged.
\end{proof}

\begin{remark}
Equation \eqref{eq:tiltgradrep} says: \emph{under the tilted measure, the gradient process
obeys the same driver-free tangent representation as in Section~\ref{sec:zero}}. All the
machinery of Sections~\ref{sec:discrete0}--\ref{sec:sobolev} therefore transfers, provided we
can sample (approximately) from $\Q$ and correct for the approximation. This is the
BSDE-side derivation of \emph{Adjoint Matching} \cite{AdjointMatching}; the
importance-weighted version with general gauge freedom is \emph{Stochastic Optimal Control
Matching} \cite{SOCM}, of which the present problem is the terminal-cost-only special case.
\end{remark}

\subsection{Discrete realization: shifted noise, detached control}

Let $v_\theta$ be the current model and $\uhat_\theta \coloneqq \sigma^\top \grad v_\theta$
the induced control. Simulate the controlled chain by \emph{shifting the noise}:
\begin{equation}\label{eq:shifted}
\widehat X_{i+1} = \widehat X_i + f(\widehat X_i,t_i)\,\Delta t_i
+ \sigma(\widehat X_i,t_i)\,\Delta\widetilde W_i,
\qquad
\Delta\widetilde W_i \coloneqq \Delta W_i + \uhat_\theta(\widehat X_i,t_i)\,\Delta t_i,
\end{equation}
with $\Delta W_i \sim \mathcal N(0,\Delta t_i I)$ the increments of the sampling
(controlled-measure) Brownian motion and $\uhat_\theta$ treated as an \emph{exogenous
constant} in the computation graph (``detached'').

\begin{proposition}[Detached backprop computes the uncontrolled tangent along tilted paths]
\label{prop:detached}
Reverse-mode differentiation of $r(\widehat X_N)$ through \eqref{eq:shifted}, with
$(\Delta W_i)$ and the values $\uhat_\theta(\widehat X_i,t_i)$ held fixed, delivers to node
$\widehat X_i$ the vector
\[
g_i = \widehat J_{N\leftarrow i}^{\,\top}\,\grad r(\widehat X_N),
\qquad
\widehat J_{i+1\leftarrow i} = I + \grad f\,\Delta t_i
+ \sum_k \grad\sigma_k\,\Delta\widetilde W^k_i ,
\]
i.e.\ the Euler discretization of the \emph{uncontrolled} variational SDE \eqref{eq:tangent}
driven by the shifted increments $\Delta\widetilde W_i$ --- precisely the discretization of
the process appearing in Theorem~\ref{thm:tiltgrad}, since
$dW = dW^\Q + Z\,dt \approx dW^{\uhat} + \uhat\,dt$. In particular, when
$\uhat_\theta = \sigma^\top\grad v$, the $g_i$ are conditionally unbiased (in the discrete
model) for $\grad v(\widehat X_i, t_i)$; no second derivatives of $v_\theta$ are ever formed.
\end{proposition}

\begin{proof}
Immediate from the chain rule applied to \eqref{eq:shifted} with $\uhat_\theta$ exogenous:
the step Jacobian is as displayed, including the cross terms
$\grad\sigma_k\,\uhat^k\Delta t$ generated by state-dependent $\sigma$ multiplying the
shifted increment. Detaching is not an approximation but the content of
Theorem~\ref{thm:tiltgrad}: the correct integrand is the tangent of the \emph{original} flow
along the tilted path. Differentiating through the control would add terms proportional to
$\grad \uhat_\theta = \grad(\sigma^\top\grad v_\theta)$, i.e.\ spurious Hessians, and the
resulting expectation would no longer equal \eqref{eq:tiltgradrep}. Equivalently, the
backward sweep solves the \emph{lean adjoint} equation
$d\tilde a_s = -\grad f(X_s,s)^\top \tilde a_s\,ds - \sum_k \grad\sigma_k(X_s,s)^\top
\tilde a_s\, d\widetilde W^k_s$ (backward, $\tilde a_1 = \grad r(X_1)$) of
\cite{AdjointMatching}.
\end{proof}

\subsection{Importance weights and their pathwise collapse}

Since we sample under $\uhat = \uhat_\theta$ rather than the exact kernel $Z$, expectations
under $\Q$ must be importance-corrected. Let $\Pmeas^{\uhat}$ denote the law of
\eqref{eq:tilted} with $a\grad v$ replaced by $\sigma\uhat$, and $W^{\uhat}$ its Brownian
motion (the sampled increments). Conditionally on $X_t = x$, combining \eqref{eq:Q} with the
Girsanov factor between $\Pmeas$ and $\Pmeas^{\uhat}$ gives
\begin{equation}\label{eq:logw}
\log w_t \;=\; r(X_1) - v(x,t)
- \int_t^1 \uhat_s^\top dW^{\uhat}_s - \tfrac12\int_t^1 \norm{\uhat_s}^2 ds ,
\qquad
\frac{d\Q}{d\Pmeas^{\uhat}}\Big|_{\F_1} = w_t .
\end{equation}

\begin{lemma}[Pathwise form of the weights]\label{lem:weights}
Along $\Pmeas^{\uhat}$-paths, with $Z_s = \sigma^\top\grad v(X_s,s)$ the exact kernel,
\begin{equation}\label{eq:residualweight}
\log w_t \;=\; \int_t^1 (Z_s - \uhat_s)^\top dW^{\uhat}_s
\;-\; \tfrac12 \int_t^1 \norm{Z_s - \uhat_s}^2\, ds .
\end{equation}
In particular $\log w_t \equiv 0$ \emph{pathwise} (not merely in expectation) iff
$\uhat_s = Z_s$ for a.e.\ $s$ along the path.
\end{lemma}

\begin{proof}
Under $\Pmeas^{\uhat}$, $dW_s = dW^{\uhat}_s + \uhat_s\,ds$, so the BSDE \eqref{eq:qbsde}
reads $dY_s = \big(Z_s^\top \uhat_s - \tfrac12\norm{Z_s}^2\big)ds + Z_s^\top dW^{\uhat}_s$.
Integrating from $t$ to $1$ with $Y_t = v(x,t)$, $Y_1 = r(X_1)$, and substituting
$r(X_1) - v(x,t)$ into \eqref{eq:logw}:
\[
\log w_t = \int_t^1\Big(Z^\top\uhat - \tfrac12\norm{Z}^2 - \tfrac12\norm{\uhat}^2\Big)ds
+ \int_t^1 (Z-\uhat)^\top dW^{\uhat}
= -\tfrac12\int_t^1\norm{Z-\uhat}^2 ds + \int_t^1 (Z-\uhat)^\top dW^{\uhat}. \qedhere
\]
\end{proof}

The weights are thus the Dol\'eans exponential of the \emph{residual} control error: as
$v_\theta \to v$ they flatten to $1$ pathwise and the effective sample size of
self-normalized importance sampling tends to $N_{\text{paths}}$. In practice one replaces the
unknown $v(x,t)$ in \eqref{eq:logw} by $v_\theta(x,t)$ --- a function of the conditioning
point only, which changes neither the conditional weighted-regression minimizer nor the
$O(1)$ scale of the weights --- self-normalizes within each time slice, and, if the effective
sample size degrades early in training, resamples: this is twisted sequential Monte Carlo
with twist $e^{v_\theta}$.

\subsection{Value targets and the Bou\'e--Dupuis bound}

\begin{proposition}[Exact bias of the controlled value target]\label{prop:elbo}
Define, along $\Pmeas^{\uhat}$-paths,
\[
V_t \coloneqq r(X_1) - \tfrac12\int_t^1 \norm{\uhat_s}^2\,ds .
\]
Then
\begin{equation}\label{eq:bias}
\E^{\uhat}\big[V_t \,\big|\, X_t\big] - v(X_t,t)
= -\tfrac12\, \E^{\uhat}\Big[\int_t^1 \norm{Z_s - \uhat_s}^2\, ds \,\Big|\, X_t\Big]
\;\le\; 0,
\end{equation}
with equality iff $\uhat = Z$: the bias is second order in the control error and one-sided.
Moreover the control-variate-corrected target
$V_t - \int_t^1 \uhat_s^\top dW^{\uhat}_s = v_\theta(X_t,t) + \log w_t^{\theta}$ (with
$v_\theta$ used in \eqref{eq:logw}) has \emph{zero variance} at the fixed point, by
Lemma~\ref{lem:weights}.
\end{proposition}

\begin{proof}
As in the proof of Lemma~\ref{lem:weights},
$r(X_1) - v(X_t,t) = \int_t^1 (Z^\top\uhat - \tfrac12\norm{Z}^2)ds + \int_t^1 Z^\top
dW^{\uhat}$. Take conditional expectations (the stochastic integral is a true martingale for
$Z\in$ BMO and admissible $\uhat$) and subtract $\tfrac12\E\int\norm{\uhat}^2$:
\[
\E^{\uhat}[V_t\mid X_t] - v(X_t,t)
= \E^{\uhat}\!\int_t^1\!\Big(Z^\top\uhat - \tfrac12\norm{Z}^2
- \tfrac12\norm{\uhat}^2\Big)ds
= -\tfrac12\,\E^{\uhat}\!\int_t^1 \norm{Z-\uhat}^2 ds .
\]
The displayed pathwise identity for the corrected target is \eqref{eq:logw} rearranged.
\end{proof}

Proposition~\ref{prop:elbo} is the conditional form of the Bou\'e--Dupuis variational formula
\cite{BoueDupuis}
\[
v(x,t) = \sup_{u\,\in\,\mathcal A}\;
\E^{u}\Big[r(X_1) - \tfrac12\int_t^1 \norm{u_s}^2\,ds \,\Big|\, X_t = x\Big],
\]
with the gradient head of the model playing the role of the (near-optimal) control; the value
bias is the Jensen gap of this evidence lower bound, and the optimizer is the $h$-transform
drift \cite{DaiPra,FlemingSoner}.

\begin{corollary}[Fixed point]\label{cor:fixed}
Consider the joint weighted Sobolev regression with targets
$\big(V_i^{\mathrm{cv}},\,g_i\big)$ from Propositions~\ref{prop:detached}
and~\ref{prop:elbo}, weights \eqref{eq:logw} (self-normalized per slice), and rollouts driven
by $\uhat_\theta = \sigma^\top\grad v_\theta$. Then $(v,\grad v)$ is a fixed point: at
$v_\theta = v$ the weights equal $1$ pathwise, the gradient targets are conditionally
unbiased for $\grad v$ by Theorem~\ref{thm:tiltgrad}, and the value targets are pathwise
exact by Proposition~\ref{prop:elbo}. (Uniqueness/convergence of the iteration is the
content of the analyses in \cite{AdjointMatching,SOCM}; the present conditioning on $X_t$,
rather than on a control-dependent initial law, removes the initial-distribution
obstruction discussed there.)
\end{corollary}

\subsection{Algorithm}

\begin{algorithm}[h]
\caption{Tangent-target regression for $v(x,t)=\log\E[e^{r(X_1)}\mid X_t=x]$}
\begin{algorithmic}[1]
\State \textbf{given} model $v_\theta$; define
$\uhat_\theta = \sigma^\top \grad_x v_\theta$ by automatic differentiation
\Loop
\State sample initial points $(x_0)$ and increments $\Delta W_i \sim \mathcal N(0,\Delta t_i I)$
\State roll out \eqref{eq:shifted} with $\uhat_\theta$ \textbf{detached};
retain the computation graph
\State $g_i \gets$ reverse-mode gradients of $r(\widehat X_N)$ at each node $\widehat X_i$
\Comment{one VJP sweep}
\State $\log w_i \gets r(\widehat X_N) - v_\theta(\widehat X_i,t_i)
- \sum_{j\ge i}\big[\uhat_j^\top \Delta W_j + \tfrac12\norm{\uhat_j}^2\Delta t_j\big]$
\State $V_i \gets r(\widehat X_N) - \sum_{j\ge i}\tfrac12\norm{\uhat_j}^2\Delta t_j
\;\big[\, -\sum_{j\ge i}\uhat_j^\top\Delta W_j \text{ as control variate}\,\big]$
\State $\bar w_i \gets \mathrm{softmax}_{\text{paths}}(\log w_i)$ per time slice
(resample if ESS low)
\State $\theta \gets \theta - \eta\,\grad_\theta \sum_i \bar w_i\Big[
\big(v_\theta(\widehat X_i,t_i)-V_i\big)^2
+ \lambda\,\norm{\grad_x v_\theta(\widehat X_i,t_i)-g_i}^2\Big]$,
targets detached
\EndLoop
\end{algorithmic}
\end{algorithm}

\begin{remark}[Practical notes]
(i) Use a single network $v_\theta$ and obtain $\grad v_\theta$ by autodiff
(\texttt{create\_graph=True}), so value and gradient heads are consistent by construction and
the learned gradient is conservative. (ii) The controlled drift $f + a\grad v$ is more
expansive than $f$ near maxima of $r$, so the tangent targets can have heavier tails than in
Section~\ref{sec:zero}; antithetic increments and the gauge (reparameterization) matrices
$M_t$ of \cite{SOCM} are the standard variance reductions. (iii) Early in training
$\uhat_\theta$ is uninformative and the weights carry the burden; a warm start with the
unweighted scheme of Section~\ref{sec:zero} applied to $u = e^{v}$ (or annealing $r$)
stabilizes the loop.
\end{remark}

\begin{thebibliography}{99}

\bibitem{Kunita}
H.~Kunita.
\newblock \emph{Stochastic Flows and Stochastic Differential Equations}.
\newblock Cambridge Studies in Advanced Mathematics 24, Cambridge University Press, 1990.

\bibitem{PardouxPeng}
E.~Pardoux and S.~Peng.
\newblock Backward stochastic differential equations and quasilinear parabolic partial
differential equations.
\newblock In \emph{Stochastic Partial Differential Equations and Their Applications},
Lecture Notes in Control and Information Sciences 176, Springer, 1992.

\bibitem{ElKarouiPengQuenez}
N.~El~Karoui, S.~Peng, and M.-C.~Quenez.
\newblock Backward stochastic differential equations in finance.
\newblock \emph{Mathematical Finance}, 7(1):1--71, 1997.

\bibitem{MaZhang}
J.~Ma and J.~Zhang.
\newblock Representation theorems for backward stochastic differential equations.
\newblock \emph{Annals of Applied Probability}, 12(4):1390--1418, 2002.

\bibitem{Kobylanski}
M.~Kobylanski.
\newblock Backward stochastic differential equations and partial differential equations with
quadratic growth.
\newblock \emph{Annals of Probability}, 28(2):558--602, 2000.

\bibitem{BriandHu}
P.~Briand and Y.~Hu.
\newblock BSDE with quadratic growth and unbounded terminal value.
\newblock \emph{Probability Theory and Related Fields}, 136:604--618, 2006.

\bibitem{Kazamaki}
N.~Kazamaki.
\newblock \emph{Continuous Exponential Martingales and BMO}.
\newblock Lecture Notes in Mathematics 1579, Springer, 1994.

\bibitem{BoueDupuis}
M.~Bou\'e and P.~Dupuis.
\newblock A variational representation for certain functionals of Brownian motion.
\newblock \emph{Annals of Probability}, 26(4):1641--1659, 1998.

\bibitem{FlemingSoner}
W.~H.~Fleming and H.~M.~Soner.
\newblock \emph{Controlled Markov Processes and Viscosity Solutions}, 2nd ed.
\newblock Springer, 2006. (Logarithmic transformations: Ch.~VI.)

\bibitem{DaiPra}
P.~Dai~Pra.
\newblock A stochastic control approach to reciprocal diffusion processes.
\newblock \emph{Applied Mathematics and Optimization}, 23:313--329, 1991.

\bibitem{ElworthyLi}
K.~D.~Elworthy and X.-M.~Li.
\newblock Formulae for the derivatives of heat semigroups.
\newblock \emph{Journal of Functional Analysis}, 125(1):252--286, 1994.

\bibitem{GriewankWalther}
A.~Griewank and A.~Walther.
\newblock \emph{Evaluating Derivatives: Principles and Techniques of Algorithmic
Differentiation}, 2nd ed.
\newblock SIAM, 2008.

\bibitem{HugeSavine}
B.~Huge and A.~Savine.
\newblock Differential machine learning.
\newblock arXiv:2005.02347, 2020.

\bibitem{OSM}
B.~Negyesi, K.~Andersson, and C.~W.~Oosterlee.
\newblock The One Step Malliavin scheme: new discretization of BSDEs implemented with deep
learning regressions.
\newblock arXiv:2110.05421; \emph{IMA Journal of Numerical Analysis}, 2024.

\bibitem{Vidales}
M.~Sabate-Vidales, D.~\v{S}i\v{s}ka, and L.~Szpruch.
\newblock Unbiased deep solvers for linear parametric PDEs.
\newblock arXiv:1810.05094; \emph{Applied Mathematical Finance}, 2021.

\bibitem{HanJentzenE}
J.~Han, A.~Jentzen, and W.~E.
\newblock Solving high-dimensional partial differential equations using deep learning.
\newblock \emph{Proceedings of the National Academy of Sciences}, 115(34):8505--8510, 2018.

\bibitem{SOCM}
C.~Domingo-Enrich, J.~Han, B.~Amos, J.~Bruna, and R.~T.~Q.~Chen.
\newblock Stochastic Optimal Control Matching.
\newblock arXiv:2312.02027, 2023.

\bibitem{AdjointMatching}
C.~Domingo-Enrich, M.~Drozdzal, B.~Karrer, and R.~T.~Q.~Chen.
\newblock Adjoint Matching: fine-tuning flow and diffusion generative models with
memoryless stochastic optimal control.
\newblock arXiv:2409.08861, 2024.

\end{thebibliography}

\end{document}
